\addtotoc{Оршил}
\chapter*{Оршил}

\section*{Судалгааны сэдвийн үндэслэл}

Орчин үеийн мэдээллийн технологи, өгөгдлийн шинжлэх ухааны (Data Science) салбарт цаг хугацааны цувааны өгөгдөл (time-series data) дээр таамаглал хийх нь судлаачдын анхаарлыг ихээр татах болсон. Үүний нэг хүндрэлтэй боловч сонирхолтой жишээ нь гадаад валютын (Forex) зах зээлийн түүхэн өгөгдөл дээр таамаглал хийх асуудал юм. Энэхүү өгөгдөл нь маш их хэмжээний шуугиантай (high noise), олон хэмжээст (high-dimensional), мөн шугаман бус (non-linear) шинж чанартай байдгаараа компьютерийн ухааны хувьд машин сургалтын (Machine learning) загваруудыг турших, оновчлох (optimization) томоохон сорилт болдог.

Уламжлалт дүн шинжилгээний аргууд нь зөвхөн санхүүгийн шинжээчдийн туршлагад тулгуурладаг бөгөөд их хэмжээний өгөгдлийг (Big Data) богино хугацаанд боловсруулж, математик статистикийн зүй тогтлыг олж илрүүлэхэд хүний нөөцийн болон тооцооллын хязгаарлагдмал байдалтай тулгардаг. Тиймээс үнийн зүй тогтлыг илрүүлэх үйл явцыг компьютерийн ухааны үүднээс "өгөгдөлд суурилсан ангиллын болон таамаглалын систем" (Data-driven classification and forecasting system) хэлбэрээр шийдвэрлэх нь орчин үеийн программ хангамжийн шийдэл болж байна.

\section*{Судалгааны шаардлага ба актуальность}

Сүүлийн арваад жилд машин сургалт (Machine Learning) болон хиймэл оюун ухааны технологиуд санхүүгийн салбарт өргөн хэрэглэгдэж эхэлсэн. Gradient Boosted Decision Trees (GBDT), нейрон сүлжээ зэрэг орчин үеийн алгоритмууд нь том хэмжээний өгөгдөл боловсруулж, нарийн төвөгтэй хэв маягуудыг илрүүлэх чадвартай болсон. Гэвч дараах асуудлууд одоо ч байсаар байна:

\begin{itemize}
    \item \textbf{Overfitting:} Олон загвар нь сургалтын өгөгдөл дээр сайн ажилладаг боловч бодит зах зээл дээр муу үр дүн үзүүлдэг.
    \item \textbf{Хүртээмж:} Ихэнх судалгаанууд академик түвшинд хийгдсэн, практик хэрэглэгчдэд хүртээмжтэй бус.
    \item \textbf{Дан загварын хязгаарлалт:} Нэг загвар ашигладаг систем нь тогтвортой бус, эрсдэл өндөртэй.
    \item \textbf{Бодит цагийн мэдээлэл дутагдал:} Загвар дангаараа хүрэлцэхгүй, хэрэглэгчдэд ойлгомжтой интерфейс, шинэ мэдээлэл хэрэгтэй.
\end{itemize}

Эдгээр асуудлыг шийдвэрлэхийн тулд ансамбль загварууд (олон загварыг нэгтгэх), walk-forward validation (цагийн дарааллаар баталгаажуулах) аргачлал, мөн гар утасны аппликейшн хөгжүүлэх нь зайлшгүй шаардлагатай болсон.

\section*{Судалгааны зорилго}

Энэхүү судалгааны ажлын гол зорилго нь \textbf{хийсвэрлэл болон өгөгдлийн шуугиан ихтэй цаг хугацааны цувааны өгөгдөл дээр ансамбль машин сургалтын (Ensemble Machine Learning) алгоритмуудыг сургах замаар хоёртын ангилал (binary classification) хийх, улмаар түүнийг бодит хугацаанд ажиллах (real-time) Backend архитектур болон хэрэглэгчийн гар утасны аппликейшнтэй (Frontend) нэгтгэсэн цогц программ хангамжийн систем хөгжүүлэх} явдал юм.

\section*{Судалгааны зорилтууд}

Дээрх зорилгод хүрэхийн тулд дараах зорилтуудыг дэвшүүлсэн:

\begin{enumerate}
    \item Компьютерийн ухааны үүднээс шуугиантай цаг хугацааны цувааны өгөгдлийг (time-series data) цуглуулж цэвэрлэх, түүнээс шинж чанар боловсруулах (Feature Engineering - техник индикатор) аргуудыг судлах.
    \item LightGBM, XGBoost, CatBoost зэрэг алгоритмуудын давуу талд үндэслэн санал өгөх механизмтай (Voting Classifier) ансамбль системийг бүтээх.
    \item 2015--2024 оны түүхэн өгөгдөл (ойролцоогоор 10 жилийн хэмжээтэй том өгөгдөл) ашиглан загваруудыг сургаж, баталгаажуулж, тестлэх туршилтууд (experimentation) хийх.
    \item Олон уулзвар бүхий цагийн дараалалтай баталгаажуулалтын (Walk-forward validation) аргачлалыг хэрэгжүүлж, загварын хэт нийцэл (overfitting) үүсэх эрсдэлийг бууруулах.
    \item Загварыг бодит цагийн өгөгдөл боловсруулах урсгал (Data pipeline) болон API (backend) бүтэцтэй холбож автоматжуулах.
    \item React Native ашиглан хэрэглэгч рүү чиглэсэн хөдөлгөөнт гар утасны (frontend) аппликейшн хөгжүүлж, серверийн талын бодит цагийн мэдээлэлгүүдийг мэдээллэх функцүүдийг нэгтгэх.
    \item Системийн ерөнхий архитектурын найдвартай байдал, тооцооллын гүйцэтгэл (latency), болон машин сургалтын загварын хувьсах хэмжигдэхүүнүүдийг үнэлэх.
\end{enumerate}

\section*{Судалгааны ач холбогдол}

Энэхүү судалгааны ажлын ач холбогдол нь:

\begin{itemize}
    \item \textbf{Онолын ач холбогдол:} Их хэмжээний шуугиантай өгөгдөл дээр ансамбль загваруудын хослол болон walk-forward validation аргачлалыг амжилттай туршиж, загварын ерөнхийлөх чадварыг (generalization) хэрхэн сайжруулсныг харуулах.
    \item \textbf{Кодчлол болон инженерийн ач холбогдол:} Бодит цагт ажиллах хүчин чадалтай (scalable backend), өгөгдөл боловсруулах урсгал, хэрэглэгчид хүртээмжтэй гар утасны аппликейшн (cross-platform frontend) зэргийг нэгтгэсэн \textbf{End-to-End} программ хангамжийн цогц шийдлийг (system architecture) боловсруулсан. 
    \item \textbf{Практик ач холбогдол:} Хүндрэлтэй тооцоолуур, санхүүгийн туршлагыг программ хангамжаар дамжуулан энгийн хэрэглэгчдэд технологийн шийдэл (FinTech) болгож өгсөн.
\end{itemize}

Судалгааны ажил нь загварын сургалтаас эхлээд мобайл аппликейшн хүртэлх бүрэн end-to-end системийг хамарсан цогц инженерийн ажил юм.
